\section{Integrazione dell'OSINT negli attacchi a dizionario}\label{sec:IntegrazioneOSINT}
\textit{Profilazione, generazione mirata con CUPP e CeWL ed efficacia dei dizionari personalizzati}
Mentre i dizionari tradizionali utilizzano elenchi generici di parole comuni, l’integrazione dell’OSINT (\cref{sec:OSINT}) permette di creare dizionari personalizzati, basati sulle informazioni personali e professionali degli utenti.

\subsection*{Utilizzo delle informazioni OSINT per costruire dizionari personalizzati}

Le informazioni raccolte con l'OSINT possono riguardare la sfera personale, come nomi di familiari, animali domestici, hobby e date significative, oppure quella professionale, come progetti, marchi, tecnologie e terminologia aziendale.

L’analisi di questi dati consente di individuare parole chiave legate al contesto dell’individuo o dell’organizzazione, utilizzabili come base per la creazione di wordlist personalizzate.
\section*{Individuazione di schemi ricorrenti nella scelta delle password}
L'OSINT fornisce gli ``ingredienti'' informativi, mentre specifiche tecniche di elaborazione applicano regole di combinazione e trasformazione per rappresentare schemi ricorrenti nella scelta delle password.
Tra le principali tecniche si possono distinguere:
\begin{enumerate}
    \item \textbf{Concatenazione e associazione di date:} combinazione di parole significative e riferimenti temporali. Gli utenti possono, ad esempio, associare elementi familiari o interessi personali a sequenze numeriche.
    \item \textbf{Sostituzione di caratteri (\textit{Leetspeak}):} trasformazione di lettere in numeri o simboli visivamente simili, secondo convenzioni ricorrenti. Tra gli esempi più comuni figurano la sostituzione della lettera \texttt{e} con \texttt{3}, della \texttt{a} con \texttt{@} e della \texttt{s} con \texttt{\$}.
    \item \textbf{Prefissi e suffissi:} aggiunta di sequenze di caratteri, numeri o simboli alle parole di partenza, spesso per soddisfare i requisiti formali delle policy di complessità delle password.
    \item \textbf{Word mangling e generazione basata su regole (\textit{rule-based generation}):} applicazione automatizzata di regole predefinite che combinano le trasformazioni precedenti, insieme alle modifiche della capitalizzazione, sulle parole di una lista iniziale.
    \item \textbf{Modellazione probabilistica:} utilizzo di modelli statistici, come i modelli di Markov, per rappresentare la probabilità di determinate sequenze di caratteri sulla base di schemi precedentemente osservati.
\end{enumerate}
È importante distinguere tra le tecniche di riconoscimento dei pattern (\textit{pattern recognition}), che identificano regolarità nei dati, e quelle di generazione basata su pattern (\textit{pattern-based generation}), che utilizzano tali regolarità per produrre nuove combinazioni.
Ad eccezione della modellazione probabilistica, che stima le probabilità a partire da password già osservate (come le grammatiche probabilistiche~\cite{weir2009} citate nell'introduzione), queste tecniche non richiedono intelligenza artificiale o machine learning: molti strumenti operano attraverso semplici regole deterministiche di trasformazione e combinazione delle stringhe.
Per automatizzare alcune di queste operazioni esistono strumenti specifici, tra cui CUPP e CeWL, che adottano approcci differenti alla costruzione di wordlist personalizzate.
\subsection*{1. CUPP (Common User Passwords Profiler)}
CUPP è uno strumento open source scritto in Python, progettato per generare wordlist personalizzate a partire da informazioni associate a un determinato profilo individuale.
Il suo funzionamento si basa principalmente su una procedura interattiva, nella quale vengono inserite informazioni quali nomi, soprannomi, date significative, interessi e altri elementi potenzialmente rilevanti. Lo strumento utilizza questi dati come parole di partenza per applicare combinazioni e trasformazioni automatiche.
Tra le tecniche implementate figurano la combinazione di parole, la variazione dei formati numerici, le modifiche della capitalizzazione e alcune sostituzioni di caratteri. Si tratta delle operazioni di \textit{word mangling} descritte sopra.
CUPP mette inoltre a disposizione funzionalità per elaborare wordlist già esistenti e generare varianti a partire da parole fornite dall'operatore.
\subsection*{2. CeWL (Custom Word List generator)}
CeWL è uno strumento open source scritto in Ruby, orientato all'estrazione di termini dai contenuti pubblicati sul web.
A differenza di CUPP, che si basa principalmente su informazioni inserite dall'operatore, CeWL analizza direttamente le pagine di un sito web, seguendo i collegamenti interni entro una profondità configurabile.
Durante la scansione, lo strumento estrae le parole presenti nei contenuti testuali delle pagine e costruisce una lista lessicale associata al sito analizzato. Possono così emergere termini specifici dell'organizzazione, come nomi di prodotti, servizi, marchi, progetti e tecnologie.
Il funzionamento e le principali opzioni da riga di comando sono riassunti in \cref{fig:cewl}.
CeWL non esegue automaticamente un'analisi semantica approfondita dei contenuti né individua necessariamente le password utilizzate dai dipendenti. La sua funzione consiste nell'estrarre termini dai contenuti web e organizzarli in una wordlist personalizzata.
L'utilità di questo approccio risiede nella possibilità di individuare il lessico caratteristico di un'organizzazione, che può essere successivamente impiegato nell'ambito di attività autorizzate di verifica della robustezza delle credenziali.
\begin{figure*}[t]
\centering
\begin{tikzpicture}[
  node distance=0.7cm,
  main/.style={draw=teal!60!black, fill=teal!10, rounded corners=3pt,
               text width=5.2cm, align=center, minimum height=1.2cm},
  start/.style={draw=gray!70, fill=gray!10, rounded corners=3pt,
                text width=5.2cm, align=center, minimum height=1.2cm},
  opt/.style={draw=violet!60!black, fill=violet!10, rounded corners=3pt,
              text width=4.2cm, align=center, minimum height=1.2cm},
  flow/.style={-{Stealth[length=2.5mm]}, thick},
  link/.style={dashed, gray}
]
\node[start] (site) {\textbf{Sito web}\\{\small Contenuti pubblici}};
\node[main, below=of site] (crawl) {\textbf{1. Scansione (crawling)}\\{\small Segue i collegamenti interni}};
\node[main, below=of crawl] (filter) {\textbf{2. Estrazione e filtro}\\{\small Selezione dei termini testuali}};
\node[main, below=of filter] (wl) {\textbf{3. Wordlist personalizzata}\\{\small Lessico specifico dell'organizzazione}};
\node[start, below=of wl] (audit) {\textbf{Verifica di sicurezza}\\{\small Analisi autorizzata delle credenziali}};
\draw[flow] (site) -- (crawl);
\draw[flow] (crawl) -- (filter);
\draw[flow] (filter) -- (wl);
\draw[flow] (wl) -- (audit);
\node[opt, right=1.4cm of site] (e) {\textbf{\texttt{-e}}\\{\small Estrazione degli indirizzi e-mail}};
\node[opt, right=1.4cm of crawl] (d) {\textbf{\texttt{-d}}\\{\small Profondità di scansione}};
\node[opt, right=1.4cm of filter] (m) {\textbf{\texttt{-m}}\\{\small Lunghezza minima delle parole}};
\node[opt, right=1.4cm of wl] (n) {\textbf{\texttt{--with-numbers}}\\{\small Inclusione di parole con cifre}};
\draw[link] (site) -- (e);
\draw[link] (crawl) -- (d);
\draw[link] (filter) -- (m);
\draw[link] (wl) -- (n);
\end{tikzpicture}
\caption{Funzionamento generale di CeWL: dalla scansione dei contenuti web alla costruzione di una wordlist personalizzata, con alcune opzioni da riga di comando.}
\label{fig:cewl}
\end{figure*}
\section*{Efficacia dei dizionari mirati}
Rispetto a una lista generica, un dizionario costruito con l'OSINT restringe le ipotesi a quelle legate alla persona o all'organizzazione analizzata, riducendo il numero di combinazioni da provare (\cref{sec:AttacchiDizionario}). Il vantaggio effettivo dipende però dalla qualità dei dati raccolti e dal metodo di verifica adottato: informazioni errate, obsolete o riferite a omonimi (\cref{sec:OSINT}) producono candidati inutili e annullano il beneficio della personalizzazione.
